%% PDF/UA-2 validation case: roughly the construct surface of tagged.tex,
%% but built as a real accessible document so the suite can run veraPDF on
%% the result instead of leaving it to a manual step before release.
%%
%% What this catches that no other case can: a structure element opened at
%% the wrong moment.  Opening the \glt language Span while still in vertical
%% mode -- rather than deferring it to \everypar, so the paragraph's own P
%% and MC are already open -- produces marked content that straddles its
%% parent.  Every geometric and flat-structure assertion in this suite
%% accepts that happily; veraPDF rejects it on all three profiles.
%%
%% Deliberately NO footnote here, although examples in footnotes are a
%% supported (and delicate) construct.  On this TeX Live a footnote in a
%% tagged UA-2 article emits content that is "neither marked as Artifact nor
%% tagged as real content" under pdflatex and xelatex -- reproducible with a
%% plain \documentclass{article} and no linguexx loaded at all, so it is a
%% kernel/latex-lab matter and not ours to assert on.  Footnote examples stay
%% covered structurally by tagged.tex and, end-to-end, by examples/ua-demo.
%% PDF/UA-2 preamble, mirroring examples/ua-demo.tex: this is the only
%% configuration whose output veraPDF is asked to validate, so it has to be
%% a realistic accessible build rather than a bare tagging switch.
%% pdfstandard=ua-2 requires PDF 2.0.  A dc:title is mandatory under
%% ISO 14289-2 clause 8.11.1, hence the pdftitle.
\DocumentMetadata{lang=en,pdfversion=2.0,pdfstandard=ua-2,tagging=on}
\documentclass{article}
\usepackage{iftex}
\ifPDFTeX
  \usepackage[T1]{fontenc}
  \usepackage[utf8]{inputenc}
\else
  \usepackage{fontspec}
\fi
\usepackage[lazy,gb4e]{linguexx}
\usepackage{hyperref}
\hypersetup{pdftitle={linguexx PDF/UA regression case},
            pdfdisplaydoctitle=true}
\pagestyle{empty}
\begin{document}

\section{UASECT accessible constructs}

UAINTRO ordinary running text precedes the first example.

\ex.\label{ua:one} UAMAIN a numbered example.
\a. UAALPHA a sub-example.
\b. *UABETA a judged sub-example.
\a. UAROMAN a sub-sub-example.
\z.\z.

%% gloss: object-language tier, a Leipzig abbreviation with its expansion,
%% and a free translation in a third language -- the \glt Span under test
\GlossTierLang{1}{de}
\GlossTransLang{en}
\exg. UAOBJ Der Hund bellte.\\ the dog bark.\lpzg{3sg.pst}\\
\glt `UATRANS The dog barked.'

%% stacked alternatives, text-mode: no Formula element may appear (that is
%% the PDF/UA-2 failure which forced the 0.12 removal of \altg)
\ex. UAALTN I saw \altn{the cat}{a dog} here.

\exg. UAALTG \altg{Socke}{Tonne} bellte.\\
      altg \altg{sock.\lpzg{sg}}{ton.\lpzg{sg}} bark.\lpzg{pst}\\
\glt `UAALTGT The sock or the ton barked.'

Reference to \ref{ua:one} resolves.

\begin{exe}
\ex[??]{UAEXE an item in a batch.}
\end{exe}

%% the abbreviation list is itself a tagged list
UALIST follows.
\lpzglist
\end{document}
